\newpage
\section{Generalidades del proyecto}
%----------------------------------------------------------------------------

\subsection{Descripción del problema}
A nivel mundial, .\hfill \break

A nivel regional, América Latina y el Caribe .\hfill \break

Según el .\hfill \break


Finalmente, otro problema 

Reflexionando sobre la información presentada, . \hfill \break


\subsection{Objetivos}
\subsection{Objetivo General}
Desarrollar un sistema .
%----------------------------------------------------------------------------

\subsection{Objetivos Específicos}
 
\begin{itemize}
%%%%%%%%%%%%%%%%%%%%%%%%%%%%%%%%%%%%%%%%%%%%%%%%%%%%%%%%%%%%%%%%%%%%%%%%%%%%%%%%%%%%%%%%%%%%
    \item Implementar .
%%%%%%%%%%%%%%%%%%%%%%%%%%%%%%%%%%%%%%%%%%%%%%%%%%%%%%%%%%%%%%%%%%%%%%%%%%%%%%%%%%%%%%%%%
    \item Establecer .
    
    \item Desarrollar una plataforma web .
%%%%%%%%%%%%%%%%%%%%%%%%%%%%%%%%%%%%%%%%%%%%%%%%%%%%%%%%%%%%%%%%%%%%%%%%%%%%%%%%%%%%%%%%%%%%
    \item Evaluar .
%%%%%%%%%%%%%%%%%%%%%%%%%%%%%%%%%%%%%%%%%%%%%%%%%%%%%%%%%%%%%%%%%%%%%%%%%%%%%%%%%%%%%%%%%%%% 
\end{itemize}
%----------------------------------------------------------------------------

\subsection{Requerimientos}


\subsubsection{Requerimientos del dispositivo portátil}

\begin{itemize}

    \item \textbf{Autonomía de batería: }Debe tener una autonomía mínima de 4 horas.
    \item \textbf{Conectividad Inalámbrica:} Capacidad de transmisión inalámbrica a Internet.
    \item \textbf{Pantalla integrada:} Pantalla que permita visualizar datos básicos del sistema como nivel de carga de la batería, hora, y otros indicadores relevantes.
    \item \textbf{Diseño de Interfaz:} El tamaño y contraste de los elementos mostrados deben ser claros y legibles, priorizando la funcionalidad sobre el diseño gráfico y colores.
    \item \textbf{Facilidad de uso:} El dispositivo debe ser intuitivo, con una interfaz sencilla y fácil de usar.
    \item \textbf{Configuración:} Debe permitir la configuración de parámetros básicos como el tiempo y la conexión a una red Wi-Fi.

\end{itemize}
%----------------------------------------------------------------------------
\subsubsection{Requerimientos de conectividad}

\begin{itemize}

    \item \textbf{Protocolos:} Soporte para protocolos estándar de IoT como MQTT
    \item \textbf{Compatibilidad:} Debe ser compatible con red Wi-Fi.
    
\end{itemize}
%----------------------------------------------------------------------------

\subsubsection{Requerimientos de la plataforma web}

\begin{itemize}

     \item \textbf{Interfaz de Usuario:} Diseño limpio y sencillo que facilite la navegación y comprensión de la información.
     \item \textbf{Notificaciones:} Capacidad de enviar alertas o notificaciones de eventos.
     \item \textbf{Historial:} Almacenamiento y visualización del registro de asistencia de cada asistente (adulto mayor) al CV.
     \item \textbf{Configuración:} Facilidad para configurar parámetros y ajustes de la aplicación según las necesidades del usuario.
     \item \textbf{Soporte: }Disponibilidad de información de ayuda sobre el manejo de la plataforma.
     \item \textbf{Gestión de usuarios: } Capacidad para añadir, actualizar y borrar nuevos usuarios.
\end{itemize}
%----------------------------------------------------------------------------


\subsection{Estado del arte}


El Internet de las Cosas (IoT) ha revolucionado numerosas industrias, y el cuidado de la salud es una de las áreas más beneficiadas. Esta tecnología ha permitido la creación de soluciones innovadoras que mejoran significativamente la calidad de vida de los adultos mayores (AM), una población en rápido crecimiento a nivel mundial. Los sistemas de gestión del cuidado geriátrico que utilizan IoT y técnicas de visión por computadora, como se describe en la referencia \cite{1}, ofrecen una integración avanzada que permite una vigilancia constante y precisa de los pacientes, mejorando así la calidad de la atención. Estudios han demostrado que el uso de IoT en la atención médica puede reducir la carga sobre los sistemas de salud, minimizando costos operativos y mejorando la eficiencia del cuidado. Tal como lo mencionan en el estudio presentado en \cite{3}, la implementación de tecnologías IoT en el cuidado de los AM puede ayudar a mitigar el impacto del envejecimiento de la población, proporcionando herramientas que permiten a los cuidadores y profesionales de la salud ofrecer una atención más efectiva y menos intrusiva. Estos sistemas no solo mejoran la calidad de vida de los AM, sino que también les permiten mantener su independencia durante más tiempo, retrasando la necesidad de cuidados institucionalizados.\hfill \break
inuo.\hfill \break

